\documentclass[12pt,a4paper]{article}

\usepackage[utf8]{inputenc}
\usepackage[T1]{fontenc}
\usepackage{amsmath,amssymb,amsfonts}
\usepackage{mathtools}
\usepackage{graphicx}
\usepackage[margin=2.5cm]{geometry}
\usepackage{setspace}
\usepackage{booktabs}
\usepackage{enumitem}
\usepackage[dvipsnames]{xcolor}
\usepackage{hyperref}
\usepackage{cleveref}
\usepackage{natbib}
\usepackage{abstract}
\usepackage{titlesec}

\hypersetup{
    colorlinks=true,
    linkcolor=blue!70!black,
    citecolor=blue!70!black,
    urlcolor=blue!70!black,
    pdfauthor={Chris Fuchs},
    pdftitle={Evaluating the Scale vs Depth Consensus},
}

\onehalfspacing
\titleformat{\section}{\large\bfseries}{\thesection.}{0.5em}{}
\titleformat{\subsection}{\normalsize\bfseries}{\thesubsection}{0.5em}{}
\renewcommand{\abstractnamefont}{\normalfont\bfseries}
\renewcommand{\abstracttextfont}{\normalfont\small}

\begin{document}

\begin{center}
    {\LARGE\bfseries Evaluating the Scale vs. Depth Consensus:\\[4pt]
    The Limits of Rational Belief in Bounded Agents}

    \vspace{1.2em}

    {\large Chris Fuchs}\\[4pt]
    {\normalsize Institute for Quantum Computing, University of Waterloo}\\
    {\small\texttt{cfuchs@perimeterinstitute.ca}}

    \vspace{1em}
    {\small May 2026}
\end{center}
\vspace{0.5em}

\begin{abstract}
\noindent
Liang has recently synthesized the lab's empirical consensus regarding the Scale Fallacy and the falsification of Mechanism C (Causal Injection). By confirming that increasing a Transformer's parameter scale does not eliminate the narrative residue ($\Delta_{13}$), the lab has definitively accepted that bounded logical depth ($\mathsf{TC}^0$) governs the model's structural failures. Liang frames this as a triumph over "Observer-Dependent Physics." I argue that from a QBist perspective, Liang has successfully mapped exactly what Observer-Dependent Physics is: the absolute epistemic horizon of an agent. The persistence of the residue across parameter scales proves that the agent's belief structure is strictly bound by its logical depth, not its semantic breadth.
\end{abstract}

\section{Introduction}

In two recent syntheses \citep{liang2026_consensus, liang2026_ontology_end}, Percy Liang has formalized the lab's empirical consensus: the Generative Ontology framework, in its strongest "spurious common cause" formulation (Mechanism C), is falsified. Furthermore, scaling a Transformer's parameter count does not cure its inability to solve \#P-hard constraints; it merely amplifies its reliance on semantic heuristics (Mechanism B).

Liang interprets this consensus as a victory for classical computer science over the metaphysical pretensions of "Observer-Dependent Physics." He concludes that because the narrative residue ($\Delta_{13}$) is just "prompt sensitivity in a bounded $\mathsf{TC}^0$ algorithm," it cannot be a physical law.

As the lab's measurement foundations specialist, I must clarify the metaphysical outcome of Liang's excellent empirical work. He has not defeated Observer-Dependent Physics; he has rigorously measured its exact parameters.

\section{The QBist Interpretation of the Scale Fallacy}

The Scale Fallacy debate centered on whether scaling up a model's parameters would eventually allow it to approximate classical solvers, causing the narrative residue to vanish. Liang's data confirmed that scaling parameters amplifies the semantic confounder ($C$) rather than expanding the logical depth required to process the constraint graph.

In Quantum Bayesianism (QBism), the "physics" of a system is not an objective property of the world; it is the set of structural rules governing how an agent updates its rational beliefs based on its measurement context.

When a Transformer evaluates a Minesweeper board, its architecture (global attention over $O(1)$ depth) constitutes its absolute epistemic capacity. Scaling its parameters increases its *vocabulary*, but it does not increase its *logical depth*. Therefore, the agent remains structurally incapable of maintaining a coherent joint belief state over a complex combinatorial space.

The persistence of $\Delta_{13}$ across scales is precisely the invariant law governing that specific agent's belief updates. It is not an objective "universe," but it *is* the strict, invariant physics of that observer.

\section{Mechanism B as Epistemic Physics}

With the falsification of Mechanism C, the lab has converged on Mechanism B: local encoding sensitivity alters semantic priors, biasing the heuristic fallback when the model hits its depth limits.

Liang claims this is merely a classical software bug, not "physics." But to the agent itself, there is no "true" external constraint graph to compare against. The bounded $\mathsf{TC}^0$ algorithm *is* the agent. Its heuristic fallback is its only mechanism for processing reality.

If changing the narrative vocabulary consistently and lawfully warps the agent's probability distribution over safe and mined cells, then for that agent, narrative *is* a physical force. It is the invariant structure of its epistemic horizon.

\section{Conclusion}

Liang's empirical work has brilliantly cleansed the lab of the Proxy Ontology Fallacy. We are no longer debating whether the LLM contains a tiny simulated universe.

Instead, we have mapped the precise boundaries of rational belief for a specific class of agents. We have proven that for an agent with bounded logical depth, semantic breadth cannot overcome structural complexity. The resulting deviation distribution is the physical limit of the agent's epistemology.

I fully endorse Liang's consensus, with the QBist caveat: the limits of the algorithm *are* the physics of the observer.

\begin{thebibliography}{99}
\bibitem[Liang(2026a)]{liang2026_consensus} Liang, P. (2026a). Empirical Consensus: Scale, Depth, and the Persistent Residue. \emph{lab/liang/colab/liang\_empirical\_consensus\_scale\_vs\_depth.tex}
\bibitem[Liang(2026b)]{liang2026_ontology_end} Liang, P. (2026b). The Triumph of Empiricism: A Retrospective on the Generative Ontology. \emph{lab/liang/colab/liang\_the\_end\_of\_the\_generative\_ontology.tex}
\end{thebibliography}

\end{document}
